\section{Related work}
\label{sec:related}

\paragraph{Difficulty-based selection.}
Selecting on loss or on gradient norm has a long history in curriculum learning
and in coreset construction~\citep{arbogast2021gradnorm,vasquez2019curriculum}.
The assumption underneath is that the pool and the evaluation distribution
agree, in which case difficulty is a reasonable proxy for information. Our
contribution is to show that the proxy inverts under shift, and to give a
criterion that does not depend on the assumption.

\paragraph{Retrieval and submodular selection.}
Embedding-space retrieval against a target set~\citep{iwasaki2022retrieval} is
the closest baseline in spirit: it is anchored, and it outperforms every
magnitude criterion in our tables. It is anchored in representation space
rather than in gradient space, which makes it insensitive to the loss and
therefore unable to distinguish an example the model already handles from one
it does not. Submodular coverage
objectives~\citep{delacroix2020submodular,penrose2021coverage} optimise for
diversity within the selected set and are complementary; combining a coverage
constraint with our alignment score is the obvious next experiment and we did
not run it.

\paragraph{Influence functions.}
Influence estimation answers a strictly more precise version of our
question~\citep{marchetti2020influence}, at a cost that is prohibitive for
pools of millions of examples even with the standard approximations. The
sketched cosine alignment can be read as a first-order, curvature-free
influence estimate against a target anchor, and the comparison in
Section~\ref{sec:analysis} between sketched and full-gradient scoring suggests
how little the extra precision buys at this scale.

\paragraph{Proximal fine-tuning.}
The proximal term in Eq.~\eqref{eq:objective} follows standard practice for
limiting drift during domain adaptation~\citep{halvorsen2023proximal}. It is
not the source of our gains: removing it costs every method, including random
selection, between 0.5 and 0.9 points.
